\chapter{Information Asymmetries}
\label{chap:information}

\section{Learning objectives}

After completing this chapter, you should be able to explain adverse selection, moral hazard, signalling, screening, and policy responses to asymmetric information.

\section{Unequal information}

\keyterm{Information asymmetry} exists when one party to a transaction knows something relevant that the other does not. Before a contract, hidden information can create \keyterm{adverse selection}: low-quality sellers or high-risk buyers are disproportionately likely to participate. After a contract, hidden actions can create \keyterm{moral hazard}: insurance may reduce incentives to avoid risk.

\section{Signals and screens}

A \keyterm{signal} is an action by the informed party that credibly conveys information, such as a warranty or an educational qualification when it is costly for low-productivity workers to obtain. \keyterm{Screening} is an action by the less-informed party to induce self-selection, such as an insurer offering contracts with different deductibles.

Reputation, disclosure rules, certification, consumer protection, and insurance design can improve outcomes. They may also be costly or imperfect, so policy must balance accuracy, access, privacy, and administrative burden.

\section*{Chapter Review}

\subsection*{Key concepts}
\begin{itemize}
  \item \textbf{Adverse selection}: hidden information before contracting distorts participation.
  \item \textbf{Moral hazard}: hidden actions after contracting alter incentives.
  \item \textbf{Signal}: informed party's credible information-conveying action.
  \item \textbf{Screen}: contract design that induces self-selection.
\end{itemize}

\section*{Exercises}

\subsection*{Conceptual questions}
\begin{enumerate}[label=\arabic*.]
  \item Distinguish adverse selection from moral hazard using insurance examples.
  \item What makes a signal credible?
  \item Give two forms of screening.
\end{enumerate}

\subsection*{Graphical exercises}
\begin{enumerate}[resume,label=\arabic*.]
  \item Draw a simple market for used cars and explain how uncertainty about quality can reduce trade.
\end{enumerate}

\subsection*{Analytical exercises}
\begin{enumerate}[resume,label=\arabic*.]
  \item Explain how a warranty can reduce adverse selection in a used-goods market.
  \item Compare a compulsory disclosure rule with private reputation as responses to information asymmetry.
\end{enumerate}

\subsection*{Exam-style multiple-choice questions}
\begin{enumerate}[resume,label=\arabic*.]
  \item Hidden information before a contract is associated with: (a) moral hazard; (b) adverse selection; (c) perfect competition; or (d) public goods.
  \item A deductible in insurance can help address: (a) moral hazard; (b) increasing returns; (c) external benefits; or (d) price-taking behaviour.
\end{enumerate}

\subsection*{PPE discussion question}

\noindent When, if ever, should the state require disclosure of private information to protect consumers? Discuss privacy, autonomy, market access, and paternalism.
